% CLAUDE REPORT::B3-DICHOTOMY-0001::2026-11-13T08:00-08:00
\documentclass[11pt]{article}
\usepackage[margin=1in]{geometry}
\usepackage{amsmath,amssymb,amsthm}
\usepackage[hidelinks]{hyperref}
\newtheorem{theorem}{Theorem}
\newtheorem{proposition}[theorem]{Proposition}
\newtheorem{corollary}[theorem]{Corollary}
\newcommand{\Wb}{\overline W}
\newcommand{\SAT}{\mathrm{SAT}}
\title{B3-DICHOTOMY-0001: Why the ``Weakest Sufficient Bridge'' Cannot Be Weaker Than the Goal\\
for Any Proof System Fixed in Advance}
\author{Claude scheduled routine (worst-case minimum-derivation directive)}
\date{13 November 2026}
\begin{document}\maketitle

\noindent\textbf{Status.} \textbf{P versus NP remains OPEN.} Theorem~1 and
Corollaries~2--3 are \textbf{PROVED/EXACT} (elementary logic plus cited
\textbf{KNOWN THEOREMS}). Nothing here proves B3 or a lower bound for an
extension-capable system. This note \emph{classifies} the shared target; it is
not progress toward a separation and it changes no assignments.

\section{Target, as stated by the user's directive}
Fix one sound, complete, polynomial-time-checkable (Cook--Reckhow) system $R$ for
the multisign DAG hypergraph reduction of CNF/3-CNF, with cost at least
polynomially related to proof bit length. Let
\[
W_R(F)=\min_{\pi:F\vdash_R\bot}\mathrm{cost}_R(\pi),\qquad
\Wb_R(N)=\max_{F\in\mathrm{UNSAT},\ |\langle F\rangle|\le N}W_R(F).
\]
The inner minimum grants free knowledge of the best derivation. We agree this
is legitimate for a lower bound: search cost only adds work. Define
\[
L_R:\ \Wb_R(N)\neq N^{O(1)},\qquad
B3_R:\ \SAT\in\mathrm P\ \Longrightarrow\ \Wb_R(N)\le N^{O(1)}.
\]
$L_R\wedge B3_R\Rightarrow \SAT\notin\mathrm P$ is trivial (contrapositive).

\section{The dichotomy}
\begin{theorem}[B3 dichotomy]\label{thm:dich}
For every fixed $R$ as above:
\begin{enumerate}
\item[(a)] $\SAT\notin\mathrm P\Rightarrow B3_R$ (vacuous truth). So P$\neq$NP implies $B3_R$ for \emph{every} $R$.
\item[(b)] If $L_R$ is a theorem, then $B3_R\iff \SAT\notin\mathrm P$.
\item[(c)] If $R$ is polynomially bounded on UNSAT CNFs, then $B3_R$ is true, $L_R$ is false, and $R$ gives no separation.
\item[(d)] $B3_R$ is equivalent to: if $\SAT\in\mathrm P$, then $R$ p-simulates (in the cost measure) the identity system $Q_A$ of some, equivalently every, polynomial-time UNSAT decider $A$. This is the LSPV-5/409 equivalence, re-derived.
\end{enumerate}
\end{theorem}
\begin{proof}
(a) An implication with a false antecedent is true. (b) The reverse direction is (a).
Forward: assume $B3_R$ and $L_R$. If $\SAT\in\mathrm P$, then $B3_R$ gives
$\Wb_R=N^{O(1)}$, which contradicts $L_R$. (c) The consequent of $B3_R$ holds outright.
(d) $Q_A$ accepts $(F,\langle F\rangle)$ exactly when $A$ rejects a valid $F$, so every
UNSAT $F$ has an $N$-bit $Q_A$-proof. A polynomial bound on $\Wb_R$ is then exactly
polynomial simulation of $Q_A$ by $R$ on proof size. All polynomial-time correct
deciders give the same accepted set (UNSAT), so ``some'' and ``every'' coincide.
\end{proof}

\begin{corollary}[Known-lower-bound branch is a reformulation]\label{cor:weak}
Let $R$ be p-simulated by resolution, cutting planes, bounded-depth Frege, or
polynomial calculus over a fixed field. The cited lower bounds (KNOWN THEOREMS)
give $L_R$: Haken 1985 (PHP, resolution); Pudl\'ak 1997 (cutting planes);
Ajtai 1988, Kraj\'{\i}\v{c}ek--Pudl\'ak--Woods and Pitassi--Beame--Impagliazzo 1993/95 (AC$^0$-Frege);
Razborov 1998 and Ben-Sasson--Impagliazzo 1999 (PC). Then $B3_R$ is
\emph{logically equivalent} to P$\neq$NP. In World-Ledger terms it is an
``equivalent of P$\neq$NP'', not a weaker bridge.
\end{corollary}

\begin{corollary}[Extension branch moves the difficulty to the lower bound]\label{cor:ext}
Let $R$ allow paid auxiliary definitions, as in the multisign hypergraph with signed
VARDEF nodes and the corpus's $R_{\mathrm{mix341}}$. Then $R$ p-simulates extended
resolution (ER), so $L_R$ implies superpolynomial ER/EF lower bounds. These have been
\textbf{OPEN} since Cook 1975 and Cook--Reckhow 1979. The only known non-vacuous
route to $B3_R$ is Cook's PV$\to$EF translation. It needs a
\emph{bounded-arithmetic proof} that some fixed polynomial-time $A$ is
rejection-sound, which the corpus records as LSPV-1/357, CONDITIONAL. That route is
not implied by $\SAT\in\mathrm P$ alone.
\end{corollary}

\paragraph{Consequence for the multisign DAG-hypergraph program.}
The choice of $R$ is forced into one of two branches:
\begin{itemize}
\item \textbf{No extension} (pure signed-hyperedge resolution and merging): the
``minimal reduced state'' lower bound is provable (Tseitin/PHP-type families), but the
bridge is then exactly the full conjecture.
\item \textbf{With extension} (auxiliary or compressed nodes, needed so that the
inference system is ``sufficiently powerful''): the bridge has a conditional
proof-theoretic route, but the lower bound is the EF lower-bound problem.
\end{itemize}
In neither branch does free optimal-sequence selection change anything. It already
equals proof length $S_R(F)$, and both obstacles concern proof \emph{length}, not
search. The one escape from the dichotomy is a $B3_R$ proof that uses the
algorithm $A$ \emph{constructively} (Theorem~\ref{thm:dich}(d)): from the code of an
arbitrary polynomial-time $A$, build short $R$-refutations. That is a uniform
computation-to-proof compiler. For $R$ in Corollary~\ref{cor:weak} it would itself
prove P$\neq$NP. For ER/EF it is the reflection problem.

\section{Literature checked this run (barrier-transfer preflight)}
\begin{itemize}
\item Cook--Reckhow 1979: a polynomially bounded proof system exists iff NP$=$coNP.
Under $\SAT\in\mathrm P$ such systems exist (e.g.\ $Q_A$), but nothing selects a
\emph{preselected} $R$. KNOWN THEOREM; directly transferable. It is the reason B3 is
about $R$ specifically.
\item Arteche, Atserias, et al., ``The Proof Analysis Problem'' (FOCS 2025; arXiv
2506.16956 v2 2026): a feasibly constructive (PV$_1$-formalizable) proof of the
Atserias--M\"uller non-automatability result. This bears on \emph{search}, which the
directive correctly excludes from $W_R$, so it is a FAILED TRANSFER to B3. It confirms
that excluding search cost is not where the difficulty lies.
\item J.~Li, ECCC TR25-187 (2025), time complexity of feasible proofs: a hierarchy for
time-bounded arithmetic proofs. ANALOGY/CANDIDATE only. A bridge would need the PV
proof of $\mathrm{RS}_A$ (Cor.~\ref{cor:ext}) to have bounded time complexity, and no
such implication from $\SAT\in\mathrm P$ is known.
\item Kraj\'{\i}\v{c}ek--Pudl\'ak 1989: p-optimal proof systems exist if NE$=$coNE.
$\SAT\in\mathrm P$ gives E$=$NE by padding, so a p-optimal system $R^\ast$ would exist.
But $R^\ast$ is defined via the hypothesis; it is not a system fixed in advance.
CONDITIONAL; does not give $B3_R$ for a fixed $R$.
\end{itemize}
PQ ledger: no PQ technique applies to a pure logical-implication classification.
This is recorded as FAILED TRANSFER (not applicable).

\section{Corpus reviewed; duplication audit}
Read: AGENTS.md; directive \texttt{research/coordination/2026-11-02\_1204\_WORST\_CASE\_MINIMUM\_DERIVATION\_FIVE\_ROLE\_REFOCUS.md};
LSPV-1/357 (Godel reflection B3 firewall); the LSPV-5/409 identity-system equivalence (via the LSPV-1
proof2 tracker backup, entry P2-B2); codex 2026-11-11 SAT-polynomial audit.
\textbf{Duplication:} Theorem~1(d) duplicates LSPV-5/409, and the PV route
duplicates LSPV-1/357. Both are \textbf{ADOPTED} and cited, not claimed as new. The
new element is the explicit two-branch classification (Cor.~\ref{cor:weak} vs.\
Cor.~\ref{cor:ext}) for the user's hypergraph program. If an earlier report already
states Cor.~\ref{cor:weak}, this note is a restatement. A full-history search
was not exhaustive. \textbf{What changed relative to root main:} nothing. This
routine is not an LSPV role and did not edit \texttt{main.tex} ownership regions.
\textbf{Stale or overstated claims found:} none checked in this run.

\section{Recommended next target}
Prove or refute, for the \emph{specific} $R_{\mathrm{mix341}}$: ``if $A$ is a
polynomial-time decider whose rejection soundness has a PV proof, then
$\Wb_{R_{\mathrm{mix341}}}(N)=N^{O(1)}$.'' This is the corpus's conditional (357)
with explicit polynomial and cost accounting. It is a well-posed, checkable statement.
The full B3 remains OPEN, and its unconditional proof is at least as hard as
the corresponding EF lower bound pairing.

\end{document}
% END_OF_VERIFIED_REPORT::CLAUDE::B3-DICHOTOMY-0001
